\chapter{Introduction}
\label{chap:introduction}

\section{Motivation}

Wilderness Search and Rescue (SAR) requires locating missing or injured people
over large and visually complex areas while time, personnel, visibility, and
accessibility are constrained. Unmanned Aerial Vehicles (UAVs) can reduce the
area that must be inspected from the ground and provide a viewpoint that is
otherwise difficult or dangerous to obtain. However, turning aerial images into a
reliable detection signal is challenging since a person may occupy
only a small fraction of the frame, vegetation and terrain create clutter, and
appearance changes substantially with illumination, altitude, season, and
weather.

Visible and thermal infrared cameras provide complementary information. The
visible stream contains texture, colour, and scene context, whereas the thermal
stream responds to emitted radiation and can retain a useful human signature
when visible contrast is weak. The WiSARD dataset was introduced specifically to support
the study of visual--thermal perception for wilderness SAR across varied
acquisition conditions~\cite{broyles2022wisard}. Complementarity, however,
does not imply that using both streams is automatically beneficial. The two
cameras are physically distinct, their fields of view and effective scales may
differ, and their projections are not perfectly registered. Moreover, either
stream may be absent or uninformative in practice. A multimodal detector must
therefore learn not only how to combine two signals, but also when and how much
to rely on each one.

These deployment concerns interact with an experimental one. The datasets and
training budgets available for UAV-based RGB--thermal detection are modest in
comparison with general-purpose detection corpora. Small changes in
initialization, batch order, stochastic modality masking, or nondeterministic
GPU kernels can then lead to visibly different final scores. If architectures
are compared through a single execution, the resulting ranking may describe a
particular training trajectory rather than a repeatable architectural effect.

\section{Problem statement}

This thesis investigates person detection from paired RGB and thermal images
on the WiSARD dataset. The main detector family is RT-DETR, adapted to a
four-channel input through two independent backbones and feature-level fusion.
The baseline sums corresponding RGB and infrared feature maps at the P3, P4,
and P5 pyramid levels. The studied alternative transforms the infrared feature
map with a Feature Alignment Module (FAM) before the same additive fusion is applied.

The FAM uses RGB and IR features jointly to predict deformable-convolution
offsets and modulation masks. This creates a spatially adaptive transformation
of the infrared representation, but the module receives no ground-truth supervision
indicating the physical geometric correspondence between the RGB and IR images.
Therefore a higher detection score is insufficient to claim that physical sensor
registration has been learned. Establishing what the module does requires inspecting
its internal variables and intervening on its levels, in addition to measuring the detector output.

The study addresses four connected forms of robustness:

\begin{enumerate}[label=\arabic*.]
  \item \textbf{modality robustness}, when RGB or IR is unavailable;
  \item \textbf{geometric robustness}, when paired streams are imperfectly
        aligned or differ in scale;
  \item \textbf{experimental robustness}, when training is repeated under
        different seeds and nondeterministic execution;
  \item \textbf{architectural robustness}, when the same fusion idea is moved
        to another detector family.
\end{enumerate}

\section{Research questions}

The thesis is organized around the following research questions.

\paragraph{RQ1 --- Experimental reliability.}
How strongly do run-to-run variability, validation design, and checkpoint
selection affect the ranking of RGB--thermal detectors, and which protocol can
support a fair architecture-development campaign?

\paragraph{RQ2 --- RT-DETR effectiveness and modality use.}
Does FAM provide a repeatable advantage over additive RT-DETR fusion, and does
the selected detector obtain a measurable benefit from the thermal stream when
all modality conditions are evaluated on the same paired frames and ground
truth?

\paragraph{RQ3 --- Mechanism and improved alignment.}
What geometric-functional transformation does FAM learn across P3, P4, and P5,
which failure modes does it exhibit, and can a reliability-conditioned local
selection between raw and aligned IR improve it more consistently than simpler
alignment, resolution, or gating interventions?

\paragraph{RQ4 --- Operational relevance.}
How do the detector's error profile, computational cost, and behavior under a
severe sensor-scale mismatch delimit its usefulness for UAV Search and Rescue?

\paragraph{RQ5 --- Architectural transferability.}
Does the FAM advantage observed in RT-DETR transfer to other detector families
under the implemented Deformable DETR, DINO, and YOLOv10 configurations?

\section{Scope and relationship with previous work}

This thesis builds on previous work on RGB--thermal detection. That
work examined DETR, Modal Dropout, the RT-DETR Additive fusion baseline, fixed
tiling, and CMX. Subsequent work introduced FAM on RT-DETR, developed a Hybrid
CMX variant, and adapted Deformable DETR to channel fusion with FAM. These
components define the starting point of the present work.

The FAM is therefore \emph{not} presented as a new architecture invented in
this thesis. It is described in full so that the document remains
self-contained, but its provenance is stated explicitly. The present
contribution is the reliability-aware validation of that module: its
performance across paired runs, its internal behavior, its failure modes, the
effect of targeted corrections, and the extent to which the result transfers
outside the original RT-DETR setting. The thesis additionally develops
Reliability-Conditioned Residual Alignment (RCRA), an extension developed in
this thesis that locally selects how much of the FAM correction to retain.

Early experiments were useful for discovering bugs and generating hypotheses, 
but many were individual runs produced before the final validation and checkpoint 
policy was defined. They are retained as an exploratory development history and are 
not mixed with the main multi-seed estimates.

\section{Contributions}

The contributions developed in the present work are:

\begin{enumerate}[label=\arabic*.]
  \item a reliability-aware experimental framework for RGB--IR detection,
        combining paired training seeds, transferred heads, training-only
        Modal Dropout, isolated processes, complete-video validation, and a
        two-stage checkpoint policy;
  \item a controlled multi-seed evaluation of six RT-DETR configurations,
        establishing the consistent internal advantage of standard FAM and its
        behavior when either modality is unavailable;
  \item a mechanistic study of FAM across pyramid levels, combining internal
        diagnostics, factorial interventions, identification of a compensated
        P5 collapse, and a five-seed bounded-offset test;
  \item a validation-led comparison of tiny-object, resolution, reliability,
        and residual-alignment hypotheses, followed by a configuration-matched
        full-data comparison of RCRA and FAM;
  \item a coordinate-controlled modality audit on the same 708 paired frames,
        distinguishing paired missing-RGB robustness from native-IR detection;
  \item an explicit delimitation of the result through paired error analysis,
        the frozen Carnation scale-mismatch stress test, a detector-only
        compute benchmark, and a controlled negative transfer result on
        dual-backbone YOLOv10.
\end{enumerate}

Exploratory DINO and Deformable DETR studies are retained as development
evidence, but are not presented as central contributions. Negative results are
integral to the final contribution: they delimit which mechanistic changes
improve internal behavior without improving detection and which RT-DETR
findings fail to transfer under the studied conditions.

\section{Thesis organization}

\Cref{chap:background} introduces the dataset, detector families, multimodal
fusion, missing-modality training, and the inherited architectural starting
point. \Cref{chap:methodology} reconstructs the subsequent exploratory
development and explains how the variability investigation led first to the
historical locked comparison and then to the two-stage development protocol.
\Cref{chap:evaluation} presents the main RT-DETR results, the validation-led
architecture campaign and modality audit, internal FAM diagnostics,
bounded-offset intervention, error analysis, Carnation stress test,
computational cost, and YOLOv10 transfer study. Finally,
\Cref{chap:discussion} answers the research questions, states the validity
boundaries, and outlines future work.
